\documentclass[11pt,respuestas,a4]{aleph-examen}
% \documentclass[11pt,a5]{aleph-examen}

% -- Paquetes extra
\usepackage{aleph-comandos}
\usepackage{systeme}


% -- Datos
\institucion{Facultad de Ciencias Exactas, Naturales y Ambientales}
\carrera{Catalogo STEM}
\asignatura{Algebra lineal}
\tema{Recuperación 1}
\autor{Andres Merino}
\fecha{Periodo 2026-1}

\logouno[0.16\textwidth]{Logos/logoPUCE_04_ac}
\definecolor{colortext}{HTML}{1957d1}
\definecolor{colordef}{HTML}{1957d1}
\fuente{montserrat}


\begin{document}

\encabezado

\section*{Indicaciones}
\begin{itemize}[leftmargin=*]
\item
	En esta actividad se evalua si el estudiante (Criterio 0.0) [enunciado del criterio].
\item
    Los ejercicios deben ser resueltos a mano, sin ayuda de resúmenes o asistentes computacionales.
\item
	Las resoluciones deben entregarse en hojas de forma física al docente.
\item
	Se encuentra prohibido el uso de cualquier otra fuente de información durante todo el examen.
\item 
    No está permitido tener ningún dispositivo electrónico encendido durante el examen, incluyendo celulares, tablets, computadoras, audífonos, relojes inteligentes, ni similares.
\item
	En caso de considerar que existe un error en la pregunta o que esta se encuentra mal planteada, se debe indicar cual es el error y justificarlo.
\item
	Todas las soluciones deben estar correctamente redactadas y explicadas.
\end{itemize}

\section*{Ejercicios}



\begin{preguntas}
\item
    Bla bla bla. \puntaje{5}

\begin{respuesta}
    Bla bla bla bla.
\end{respuesta}


\end{preguntas}

\end{document}
